\documentclass[11pt,a4paper]{article}

\usepackage[margin=2.5cm]{geometry}
\usepackage[T1]{fontenc}
\usepackage[utf8]{inputenc}
\usepackage{lmodern}
\usepackage{microtype}
\usepackage{enumitem}
\usepackage{xurl}
\usepackage[hidelinks]{hyperref}
\usepackage{parskip}
\usepackage{xcolor}
\usepackage{titlesec}

\setcounter{secnumdepth}{3}
\setcounter{tocdepth}{3}
\setlist[itemize]{leftmargin=2em}

% Consistent visual styling for the document and the file catalogue.
\definecolor{brandblue}{HTML}{145A86}
\definecolor{lightblue}{HTML}{EAF3F8}
\definecolor{ruleblue}{HTML}{A9C9DA}
\definecolor{mutedtext}{HTML}{4F6470}
\titleformat{\section}
  {\Large\bfseries\color{brandblue}}
  {\thesection}{0.7em}{}
\titlespacing*{\section}{0pt}{2.2em}{0.8em}
\newcommand{\filelabel}[1]{\textbf{\color{brandblue}#1}}

\title{Asian Transport Database Overview}
\author{Summer 2026 Internship -- Darcey Sutcliffe}
\date{}

\begin{document}

\maketitle
\tableofcontents
\newpage


\section{Literature review}

Completed a literature review to understand datasets/ complete work already published and open source.\par

Outputs:\par

Asia Transport Literature Tracker.xlsx\par

\begin{itemize}
\item Includes relevant publications, reports, datasets, organisations, list of countries in study area and modes relevant to each country.
\end{itemize}

Asia Transport Literature Tracker.pdf\par

\begin{itemize}
\item Summary of spreadsheet
\end{itemize}



\section{Airports}

\subsection{Data sources used}

\subsection{Analysis done}

Airport node spatial data correction (completed in Airport.py)\par

The World Bank airport dataset provided the only identified open source data containing flight connection paths between airports. Therefore, the World Bank airport points were used as the basis for the airport nodes used for the flight network.\par

However, the World Bank airport coordinates were found to have poor spatial accuracy. The airport locations were therefore checked and corrected using the OurAirports dataset, which appeared to provide better spatial accuracy.\par

The analysis used airport data from Ourairports.csv, worldbank\_filtered\_airport\_flows.gpkg and worldbank\_filtered\_airport\_volume.gpkg.\par

OurAirports data was first filtered to include only countries within the study area using Countries\_list.xlsx. The World Bank airport coordinates were then corrected using the OurAirports coordinates where a matching IATA code was available. Where no matching OurAirports record was available, the original World Bank coordinates were retained.\par

The origin and destination coordinates of the World Bank flight connections were then corrected using the corrected airport coordinates.\par

airport\_coordinate\_audit.xlsx contains the audit of this matching process. The match\_summary sheet provides a summary of the number of airport locations that were common between the two datasets. 14 airports in the World Bank dataset could not be matched and therefore could not be spatially corrected using OurAirports.\par

\subsubsection{Shapefile matching (OSM\_extract\_country\_loop and airport\_shapefile\_nearest.py):}

Relevant tags were extracted from OSM.\par

For Aiports these were:\par

\{"aeroway": ["aerodrome"]\},\par

\{"aeroway": ["terminal"]\},\par

\{"aeroway": ["runway"]\},\par

\{"aeroway": ["taxiway"]\},\par

From here, any extracted OSM features located within 2 km of an airport node in worldbank\_filtered\_airport\_volume.gpkg were assigned to that airport. The airport’s IATA code was then added to the feature’s attributes. The final output combines all extracted OSM airport features into a single layer, with attributes identifying the feature’s original OSM tag (e.g. aerodrome, terminal, runway or taxiway) and the airport it was assigned to.\par



\subsubsection{Still to do:}

\begin{itemize}
\item Check the spatial location of the 14 airport nodes in airport\_flows\_world\_bank\_corrected.gpkg, which did not have a match to the OurAirports database. (Validating that the flows also update if any airports get updated)
\end{itemize}

\begin{itemize}
\item Add missing polygon data to the XX airports
\end{itemize}

\section{Maritime}

\subsection{Data sources used/outputted}

\begin{itemize}
\item Port watch
\end{itemize}

\begin{itemize}
\item Multi-hazard risk to global port infrastructure and resulting trade and logistics losses, [Jasper Verschuur](\url{https://www-nature-com.ezproxy-prd.bodleian.ox.ac.uk/articles/s43247-022-00656-7\#auth-Jasper-Verschuur-Aff1})
\end{itemize}

\begin{itemize}
\item Port \_landuse.gpkg ( shapefiles for ports)
\end{itemize}

\begin{itemize}
\item Global\_maritime\_network.gpkg (each node is at the centre of the port\_landuse shapefiles that it is relevant to)
\end{itemize}

\begin{itemize}
\item Ports.shp
\end{itemize}

\subsection{Analysis done}

\subsubsection{Maritime network (completed in ports.py):}

Code initally copied from the African-transport-dataset/scripts/preprocess/ "ports\_new\_merge.py". Using the Ports.shp and updated 2023-2025 attribute data provided by Jasper Verschuur (port\_calls\_average\_2023-2025.csv, port\_calls\_average\_2023-2025.csv, port\_calls\_average\_2023-2025.csv). A nodel and edge network was created, with the attributes. The network was broken up into 3 clusters, to aid visuals (northern\_russia, pacific\_network and large\_cluster). The output contains two node types, ports (land) and maritime (networking), alongside the edge connections between.\par

-file output names\par

global\_maritime\_network\_PROVA\_NEW1\par

asia \& pacific\_maritime\_network\_PROVA\_NEW1\par

northern\_russia\_network\_maritime\_network\par

large\_cluster\_maritime\_network\par

pacific\_network\_maritime\_network\par



\subsubsection{Shapefile correction (all done using port\_shapefile\_nearest.py and manually):}

port\_landuse.gpkg shapefiles were matched to their corresponding ports by assigning the nearest port node to each polygon (nodes stored in asia\_pacific\_maritime\_network\_PROVA\_NEW1.gpkg) . The results of this automated matching are stored in port\_landuse\_nearest\_port\_id\_lookup.xlsx.\par



The accuracy of these matches was then manually checked and recorded in port\_landuse\_nearest\_port\_id\_lookup\_checked.xlsx. Each match was assessed using the distance between the port node and the land-use polygon, the names associated with both datasets, and the port's location on Google Maps. In the checked file, the checked column indicates whether the port land-use shapefile (Column A) has been matched to the correct port node (Column H). Where checked indicates that the match is incorrect, the correct port id column records the correct port ID, identified manually.\par



The corrections were then applied, and the port\_landuse\_nearest\_port\_ids.gpkg and port\_landuse\_nearest\_port\_id\_lookup\_checked\_result.xlsx contain the correct results.\par



14 ports from the asia\_pacific\_maritime\_network\_PROVA\_NEW1.gpkg file did not have any polygon data (from port\_landuse.gpkg) matched to it. Seen in port\_landuse\_nearest\_port\_id\_lookup\_checked\_result.xlsx, sheet-node\_polygon\_audit, column missing\_polygon\_data = TRUE.\par



\subsubsection{Still to do:}

\begin{itemize}
\item Add missing polygon data to the 14 ports.
\end{itemize}



\section{Road}

\section{Rail}

\section{Inland waterways}

Investigated OSM tags, best appeared to be:\par

Tag = boat:yes\par

Tag = ferry : terminal\par



\section{Multimodal}



\section{ARC}

ARC is separated into two file directories:\par

\begin{quote}\ttfamily cd /data/cenv-opsis-astdb/\textbackslash\{\}<USERNAME>/asia-transport-database\end{quote}

\begin{quote}\ttfamily cd /data/cenv-opsis-astdb/\textbackslash\{\}<USERNAME>/open-gira\end{quote}



\section{Using open-gira on ARC}

Purpose: Used to create the road and rail networks using OpenGIRA.\par

\medskip
1. Go to the OpenGIRA directory:

\begin{quote}\ttfamily cd /data/cenv-opsis-astdb/\textbackslash\{\}<USERNAME>/open-gira\end{quote}

\medskip
2. Activate the OpenGIRA environment:

\begin{quote}\ttfamily module load Mamba/23.11.0-0\end{quote}

\begin{quote}\ttfamily export CONDA\_ENVS\_PATH=/home/mans4968/.mamba/envs\end{quote}

\begin{quote}\ttfamily export MAMBA\_ROOT\_PREFIX=/home/mans4968/.mamba\end{quote}

\begin{quote}\ttfamily eval "\$(conda shell.bash hook)"\end{quote}

\begin{quote}\ttfamily conda activate /home/mans4968/.mamba/envs/open-gira\end{quote}

\medskip
3. Run OpenGIRA using Snakemake and the ARC profile (not run on a .sh file):

\begin{quote}\ttfamily snakemake --profile config/ARC --cores 1\end{quote}

A specific output can also be requested, for example:\par

\begin{quote}\ttfamily snakemake --profile config/ARC --cores 1 results/russia-latest\_filter-road-primary/edges.gpq\end{quote}

\medskip
4. Check jobs currently running:

\begin{quote}\ttfamily squeue -u \$USER\end{quote}

\medskip
5. Check completed/failed jobs:

\begin{quote}\ttfamily sacct -u \$USER --format=JobID,JobName,State,Elapsed,ExitCode\end{quote}

\medskip
6. Outputs are stored in:

\begin{quote}\ttfamily /data/cenv-opsis-astdb/\textbackslash\{\}<USERNAME>/open-gira/results/\end{quote}

Completed final networks are organised into:\par

\begin{quote}\ttfamily results/road/\end{quote}

\begin{quote}\ttfamily results/rail/\end{quote}

Other folders such as geoparquet/, slices/, json/ and input/\par

contain input or intermediate OpenGIRA data.\par

 \par

\section{Asia-transport-database on ARC}

Purpose: Used to run the Asia transport database processing, including the\par

country-by-country OSM infrastructure extraction.\par

1. Go to the Asia database:\par

\begin{quote}\ttfamily cd /data/cenv-opsis-astdb/\textbackslash\{\}<USERNAME>/asia-transport-database\end{quote}

2. Input OSM data are stored in:\par

   incoming\_data/osm/\par

3. Run the required Asia transport database script/job on ARC.\par

See various .sh files\par

4. Processed outputs are stored in:\par

   processed\_data/infrastructure/\par



\section{Already existing country multimodal transport packages}

Multimodal transport data packages have already been extracted for China, Singapore, New Zealand, and Australia. These datasets can be used later for comparison and validation against the wider transport database. The packages are stored under:\par

incoming\_data/infrastructure/[folder name]\par

The content of each country package varies, and not all packages contain data for every transport mode.\par



\section{Files Overview}

\begin{center}
\colorbox{lightblue}{\parbox{0.93\textwidth}{
\centering
{\Large\bfseries\color{brandblue} Asia-transport-database repository}\\[0.35em]
{\small\color{mutedtext} A catalogue of inputs, intermediate files, final outputs and processing scripts}
}}
\end{center}

\vspace{0.4em}
\noindent\textit{Legend:} entries marked with a blue star (\textcolor{brandblue}{\textbf{★}}) are final outputs.\\[-0.2em]
\noindent\textcolor{mutedtext}{Files are grouped by their repository location; the descriptions below record their purpose, provenance and producing script.}

% Turn the existing entry breaks into quiet catalogue rules without changing
% the content of the overview.
\begingroup
\renewcommand{\medskip}{\par\vspace{0.65em}\noindent\textcolor{ruleblue}{\rule{\textwidth}{0.45pt}}\par\vspace{0.15em}}
\small



\medskip
1. \textbackslash\{\}*airport\_flows\_world\_bank\_corrected.gpkg
Directory: processed\_data/infrastructure/airport/\par

Description:\par
World Bank airport flow data with origin and destination coordinates corrected using the corrected airport locations.\par

Reference:\par
Derived from World Bank and OurAirports.\par

Produced by:\par
Airports.py\par

\medskip
2. \textbackslash\{\}*airport\_osm\_features\_within\_2km.gpkg
Directory: processed\_data/infrastructure/airport/\par

Description:\par
OSM airport lines and polygons located within 2 km of corrected World Bank airport points and assigned to the relevant airport.\par

Reference:\par
Derived from OpenStreetMap and corrected World Bank airport data.\par

Produced by:\par
airport\_shapefile\_nearest.py\par

\medskip
3. airport\_osm\_features\_within\_2km.xlsx
Directory: processed\_data/infrastructure/airport/\par

Description:\par
Audit workbook containing matched OSM features, airport summaries and OSM tag summaries.\par

Reference:\par
Derived from OpenStreetMap and corrected World Bank airport data.\par

Produced by:\par
airport\_shapefile\_nearest.py\par

\medskip
4. airport\_coordinate\_audit.xlsx
Directory: processed\_data/infrastructure/airport/\par

Description:\par
Audit spreadsheet showing which World Bank airport locations could be matched to the OurAirports dataset and therefore spatially corrected. It also allows unmatched World Bank airport points to be identified.\par

Reference:\par
Derived from World Bank and OurAirports.\par

Produced by:\par
Airports.py\par

\medskip
5. [country/region]\_[tag].parquet
Directory: processed\_data/infrastructure/osm\_filter/\par

Description:\par
OSM features extracted for specified tags from the country or regional PBF files. Outputs are produced for individual countries and, where specified, for the wider study area.\par

Reference:\par
Derived from Geofabrik OpenStreetMap extracts.\par

Produced by:\par
OSM\_extract\_country\_loop.py\par

\medskip
6. [country/region].osm.pbf
Directory: incoming\_data/osm/\par

Description:\par
Raw country or regional OpenStreetMap PBF extracts used as the source data for OSM feature extraction.\par

Reference:\par
\href{https://download.geofabrik.de/index.html}{\url{https://download.geofabrik.de/index.html}}\par

Source/Produced by:\par
Geofabrik / OpenStreetMap\par

\medskip
7. Countries\_list.xlsx
Directory: incoming\_data/\par

Description:\par
List of all countries included in the study area.\par

Reference:\par
Created by Darcey Sutcliffe.\par

Source/Produced by:\par
Created by Darcey Sutcliffe\par

\medskip
8. filtered\_Ourairports.csv
Directory: processed\_data/infrastructure/airport/\par

Description:\par
OurAirports data filtered to include small, medium and large airports within the study area.\par

Reference:\par
Derived from OurAirports.\par

Produced by:\par
Airports.py\par

\medskip
9. filtered\_Ourairports.gpkg
Directory: processed\_data/infrastructure/airport/\par

Description:\par
Spatial version of the filtered OurAirports data containing small, medium and large airports within the study area.\par

Reference:\par
Derived from OurAirports.\par

Produced by:\par
Airports.py\par

\medskip
10. global\_maritime\_network\_PROVA\_NEW1.gpkg
Directory: processed\_data/infrastructure/port/\par

Description:\par
Global maritime network containing port and maritime network nodes and the connections between them.\par

Reference:\par
Derived from maritime port data and PortWatch attributes.\par

Produced by:\par
ports.py\par

\medskip
11. \textbackslash\{\}*cleaned\_asia\_pacific\_maritime\_network\_PROVA\_NEW1.gpkg
Directory: processed\_data/infrastructure/port/\par

Description:\par
Final cleaned Asia-Pacific maritime network created by combining the three separate network clusters.\par

Reference:\par
Derived from the northern Russia, large cluster and Pacific maritime networks.\par

Produced by:\par
ports.py\par

\medskip
12. asia\_pacific\_maritime\_network\_PROVA\_NEW1.gpkg
Directory: processed\_data/infrastructure/port/\par

Description:\par
Asia-Pacific maritime network covering the study area and including connections between the separate network clusters.\par

Reference:\par
Derived from maritime port data and PortWatch attributes.\par

Produced by:\par
ports.py\par

\medskip
13. large\_cluster\_maritime\_network.gpkg
Directory: processed\_data/infrastructure/port/\par

Description:\par
Maritime network representing the large cluster containing the majority of the Asian network nodes.\par

Reference:\par
Derived from maritime port data and PortWatch attributes.\par

Produced by:\par
ports.py\par

\medskip
14. northern\_russia\_network\_maritime\_network.gpkg
Directory: processed\_data/infrastructure/port/\par

Description:\par
Maritime network representing the northern Russia cluster.\par

Reference:\par
Derived from maritime port data and PortWatch attributes.\par

Produced by:\par
ports.py\par

\medskip
15. pacific\_network\_maritime\_network.gpkg
Directory: processed\_data/infrastructure/port/\par

Description:\par
Maritime network representing the Pacific cluster.\par

Reference:\par
Derived from maritime port data and PortWatch attributes.\par

Produced by:\par
ports.py\par

\medskip
16. port\_calls\_average\_2023-2025.csv
Directory: incoming\_data/Global port supply-chains/\par

Description:\par
Port attribute data containing average port calls for 2023–2025.\par

Reference:\par
PortWatch, provided by Jasper Verschuur.\par
\href{https://portwatch.imf.org/}{\url{https://portwatch.imf.org/}}\par

\medskip
17. port\_capacity\_called\_average\_2023-2025.csv
Directory: incoming\_data/Global port supply-chains/\par

Description:\par
Port attribute data containing average port capacity called for 2023–2025.\par

Reference:\par
PortWatch, provided by Jasper Verschuur.\par
\href{https://portwatch.imf.org/}{\url{https://portwatch.imf.org/}}\par

\medskip
18. port\_landuse.gpkg
Directory: incoming\_data/infrastructure/port/\par

Description:\par
Polygon data representing maritime port land use.\par

Reference:\par
\href{https://data.mendeley.com/datasets/kdyt24tsh5/1}{\url{https://data.mendeley.com/datasets/kdyt24tsh5/1}}\par

Source:\par
Jasper Verschuur\par

\medskip
19. port\_landuse\_nearest\_port\_id\_lookup.xlsx
Directory: processed\_data/infrastructure/port/\par

Description:\par
Lookup showing the nearest port assigned to each port land-use polygon before manual corrections were applied.\par

Reference:\par
Derived from port land-use data and the maritime network.\par

Produced by:\par
port\_shapefile\_nearest.py\par

\medskip
20. port\_landuse\_nearest\_port\_id\_lookup\_checked.xlsx
Directory: incoming\_data/infrastructure/port/\par

Description:\par
Manually checked spreadsheet assessing whether the nearest port assigned to each port land-use polygon was correct. Where an incorrect match was identified, the correct port\_id was recorded.\par

Reference:\par
Created by Darcey Sutcliffe.\par

Source/Produced by:\par
Manual validation by Darcey Sutcliffe\par

\medskip
21. port\_landuse\_nearest\_port\_id\_lookup\_checked\_result.xlsx
Directory: processed\_data/infrastructure/port/\par

Description:\par
Validation spreadsheet showing the original and final port\_id for each port land-use polygon, allowing confirmation that the manual corrections were applied correctly.\par

Reference:\par
Derived from the manually checked port lookup.\par

Produced by:\par
port\_shapefile\_nearest.py\par

\medskip
22. \textbackslash\{\}*port\_landuse\_nearest\_port\_ids.gpkg
Directory: processed\_data/infrastructure/port/\par

Description:\par
Final GeoPackage containing the corrected relationship between each port land-use polygon and its corresponding port\_id.\par

Reference:\par
Derived from port land-use data and the maritime network, with manual corrections applied.\par

Produced by:\par
port\_shapefile\_nearest.py\par

\medskip
23. port\_turn\_around\_time\_average\_2023-2025.csv
Directory: incoming\_data/Global port supply-chains/\par

Description:\par
Port attribute data containing average port turnaround time for 2023–2025.\par

Reference:\par
PortWatch, provided by Jasper Verschuur.\par
\href{https://portwatch.imf.org/}{\url{https://portwatch.imf.org/}}\par

\medskip
24. Ports.csv
Directory: incoming\_data/infrastructure/port/\par

Description:\par
Point locations of maritime ports.\par

Reference:\par
\href{https://data.mendeley.com/datasets/kdyt24tsh5/1}{\url{https://data.mendeley.com/datasets/kdyt24tsh5/1}}\par

Source:\par
Jasper Verschuur\par

\medskip
25. Ourairports.csv
Directory: incoming\_data/infrastructure/Airport/\par

Description:\par
OurAirports airport attributes and coordinates. Export of all global data reported by OurAirports.\par

Reference:\par
\href{https://ourairports.com/data/}{\url{https://ourairports.com/data/}}\par

Source/Produced by:\par
External dataset\par

\medskip
26. utils\_new.py
Directory: scripts/preprocess/ (if applicable)\par

Description:\par
Utility functions used by the preprocessing scripts. The script was copied from the African-transport-dataset project.\par

Reference:\par
African-transport-dataset/scripts/preprocess/\par

Produced by:\par
Copied and adapted from the African-transport-dataset project.\par

\medskip
27. worldbank\_filtered\_airport\_flows.gpkg
Directory: incoming\_data/infrastructure/Airport/\par

Description:\par
World Bank airport flow data filtered to include only flows between airports within the study area.\par

Reference:\par
\href{https://datacatalog.worldbank.org/search/dataset/0038117/global-airports}{\url{https://datacatalog.worldbank.org/search/dataset/0038117/global-airports}}\par

Source/Produced by:\par
External dataset\par

\medskip
28. worldbank\_filtered\_airport\_volume.gpkg
Directory: incoming\_data/infrastructure/Airport/\par

Description:\par
World Bank airport data filtered to include only countries within the study area.\par

Reference:\par
\href{https://datacatalog.worldbank.org/search/dataset/0038117/global-airports}{\url{https://datacatalog.worldbank.org/search/dataset/0038117/global-airports}}\par

Source/Produced by:\par
External dataset\par

\medskip
29. \textbackslash\{\}*world\_bank\_airports\_corrected.gpkg
Directory: processed\_data/infrastructure/airport/\par

Description:\par
World Bank airport locations corrected using OurAirports coordinates where a matching IATA code was available. Original World Bank coordinates are retained where no matching OurAirports record was found.\par

Reference:\par
Derived from World Bank and OurAirports.\par

Produced by:\par
Airports.py\par

\endgroup

\end{document}
